\subsection{有限生成扩张}

命题17. --- 设 E 是域 K 的一个扩域，B 是 E 在 K 上的一个超越基。为使 E 在 K 上是有限生成的（V，第11页，定义2），充要条件是 B 有限且次数 {[}E : K(B){]} 有限。

设 E 在 K 上是有限生成的，并设 S 是 E 的一个有限子集，使得 \(\mathrm{E} = \mathrm{K}(\mathrm{S})\) 。由定理2（V，第109页），S 包含 E 在 K 上的一个超越基 \(B'\) ，且它与 B 有相同的基数（V，第110页，定理3）。因此 B 是有限的。设 \(\mathrm{K}' = \mathrm{K}(\mathrm{B})\) ；则 E 在 \(K'\) 上是代数的，且我们有 \(\mathrm{E} = \mathrm{K}'(\mathrm{S})\) ；由于 S 有限，V 第18页的定理2表明 \([E : K']\) 是有限的。

反之，设 B 有限且 \([E:K(B)]\) 有限。若 C 是向量空间 E 在 \(K(B)\) 上的一个（有限）基，则我们有 \(E = K(B)(C) = K(B \cup C)\) ，于是 E 是 K 的一个有限生成扩张。

推论1. --- 设 E 是 K 的一个有限生成扩张，并设 K' 是 K 在 E 中的相对代数闭包（V，第19页）。则 K' 在 K 上是有限次的。

设 B 是 E 在 K 上的一个超越基。由 V 第114页的推论2， \(K'\) 与 \(\mathrm{K}(\mathrm{B})\) 在 K 上线性不相交，因此 \([\mathrm{K}' : \mathrm{K}] = [\mathrm{K}'(\mathrm{B}) : \mathrm{K}(\mathrm{B})] \leq [\mathrm{E} : \mathrm{K}(\mathrm{B})]\) ，于是 \([\mathrm{K}' : \mathrm{K}]\) 的有限性由 \([\mathrm{E} : \mathrm{K}(\mathrm{B})]\) 的有限性得出。

推论2. --- 一个在其素子域上有限生成的域只包含有限多个单位根。

由推论1，我们归结为证明：一个域 L，若它是其素子域的有限次扩张，则只含有有限多个单位根。当 L 的特征 \(\neq0\) 时这是显然的，因为此时它是有限的。若 L 的特征为 0 且包含无穷多个单位根，则它包含任意高阶的本原单位根。由 V 第84页定理2，因此存在无穷多个整数 n \textgreater{} 0 使得 \(\varphi(n) \leqslant [L : Q]\) ，这是荒谬的（V，第80页，公式（2）和（3））。

推论3. --- 若 E 是域 K 的一个有限生成扩张，则 E 的每个子扩张 E' 都是有限生成的。

设 \(B'\) 是 \(E'\) 在 \(K\) 上的一个超越基。由 \(V\) 第109页定理2， \(B'\) 包含于一个超越基 \(B\) （ \(E\) 在 \(K\) 上的）中，因而由命题17是有限的。由于 \(E'\) 在 \(K(B')\) 上是代数的，且 \(E\) 是 \(K(B')\) 的一个有限生成扩张，推论1表明 \([E': K(B')]\) 有限。现在命题17表明 \(E'\) 在 \(K\) 上是有限生成的。

命题17可以用如下方式转述：K 的一个有限生成扩张是纯超越扩张 \(\mathrm{K}(x_{1},\ldots,x_{n})\) 的一个有限次代数扩张。

